\section{Core Verified Properties}

The fundamental properties of the cycle integral that hold for every
position.

\begin{itemize}
  \item Next position: $CI_{i+1} = CI_i + \operatorname{Cycle}(L)_{i+1}$ ---
    Subsection~4.1
  \item Full cycle shift: adding one cycle period advances by the total sum
    --- Subsection~4.2
  \item Sum of mod values: the modulo definition matches the list sum ---
    Subsection~4.3
  \item Strictly increasing: positive base values force a strictly growing
    integral --- Subsection~4.4
  \item Positivity: a non-negative start and positive base values keep the
    integral positive everywhere --- Subsection~4.5
  \item Unit-cycle generation: the unit cycle $[1]$ enumerates the
    consecutive integers, strictly increasing --- Subsection~4.6
\end{itemize}

\subsection{Next Position}

For any positive position, the cycle integral at that position equals the
previous value plus the current cycle element.

\begin{equation*}
\forall\, i > 0: \ \mathit{CI}(L, init)_i = \mathit{CI}(L, init)_{i-1} + \operatorname{Cycle}(L)_i
\end{equation*}

This follows directly from the recursive definition of
\texttt{CycleIntegral}.

{\raggedright
This property is verified in the
\href{https://github.com/thiagomata/prime-numbers/blob/master/src/main/scala/v1/chapter4/cycle/integral/recursive/properties/CycleIntegralProperties.scala}{\texttt{CycleIntegralProperties::\allowbreak assertNextPosition}}.
\par}

\subsection{Same Difference After Full Cycle}

The difference between consecutive values is invariant under adding a full
cycle size to both positions.

\begin{equation*}
\forall\, i \geq 0: \ \mathit{CI}(L, init)_{i+1} - \mathit{CI}(L, init)_i = \mathit{CI}(L, init)_{i+size+1} - \mathit{CI}(L, init)_{i+size}
\end{equation*}

{\raggedright
This property is verified in the
\href{https://github.com/thiagomata/prime-numbers/blob/master/src/main/scala/v1/chapter4/cycle/integral/recursive/properties/CycleIntegralProperties.scala}{\texttt{CycleIntegralProperties::\allowbreak assertSameDiffAfterCycle}}.
A key Scala verification excerpt is in Appendix~A.6; the complete proof is
linked in the source reference.
\par}

\subsection{Sum of Mod Values as List}

The cycle integral at any position equals the sum of a constructed list
containing the initial value and all cycle values up to that position (using
modular indexing for positions beyond one cycle).

\begin{equation*}
\begin{aligned}
\forall\, i \geq 0: \ \mathit{CI}(L, init)_i &= \operatorname{sum}([init] \\
&\quad + [\operatorname{Cycle}(L)_0, \operatorname{Cycle}(L)_1, \dots, \operatorname{Cycle}(L)_i])
\end{aligned}
\end{equation*}

{\raggedright
This property is verified in the
\href{https://github.com/thiagomata/prime-numbers/blob/master/src/main/scala/v1/chapter4/cycle/integral/recursive/properties/CycleIntegralProperties.scala}{\texttt{CycleIntegralProperties::\allowbreak assertSumModValueAsListEqualsCycleIntegralLoop}}.
A key Scala verification excerpt is in Appendix~A.7; the complete proof is
linked in the source reference.
\par}

\subsection{Cycle Integral Strictly Increasing}

When the initial value is non-negative and every base-list value is
positive, the cycle integral is strictly increasing: a later position
always produces a larger value.

\begin{equation*}
\begin{aligned}
init \geq 0 \;\land\; (\forall x \in L,\ x > 0) \;\land\; b > a \implies
\operatorname{CycleIntegral}(L, init)_b > \operatorname{CycleIntegral}(L, init)_a
\end{aligned}
\end{equation*}

\textbf{Proof.} Induct on $b-a$. The base case is one positive step; the
inductive step adds one more positive cycle value.

\begin{equation*}
\begin{aligned}
b=a+1 &\implies CI_b-CI_a=\operatorname{Cycle}(L)_b>0 \\
      &\quad \text{[\S3.1 and cycle positivity]} \\
       &\implies CI_b>CI_a, \\
CI_{b-1}>CI_a,\quad CI_b-CI_{b-1}=\operatorname{Cycle}(L)_b>0
       &\implies CI_b>CI_{b-1}>CI_a \\
\therefore\ CI_b &> CI_a.
  \quad \blacksquare\ \text{[Q.E.D.]}
\end{aligned}
\end{equation*}

{\raggedright
This property is verified in the
\href{https://github.com/thiagomata/prime-numbers/blob/master/src/main/scala/v1/chapter4/cycle/integral/recursive/properties/CycleIntegralProperties.scala}{\texttt{CycleIntegralProperties::\allowbreak assertCycleIntegralIncreasing}}.
\par}

\subsection{Cycle Integral Positivity}

When the initial value is non-negative and every base-list value is
positive, the cycle integral is positive at every position.

\begin{equation*}
\begin{aligned}
init \geq 0 \;\land\; (\forall x \in L,\ x > 0) \implies
\operatorname{CycleIntegral}(L, init)_i > 0
\end{aligned}
\end{equation*}

\textbf{Proof.} At the first position, the non-negative initial value and a
positive cycle value give a positive integral. Every later position adds one
more positive cycle value.

\begin{equation*}
\begin{aligned}
CI_0 &= init+\operatorname{Cycle}(L)_0>0, \\
CI_{i-1}>0,\quad CI_i-CI_{i-1}=\operatorname{Cycle}(L)_i>0
  &\implies CI_i>0 \\
\therefore\ CI_i &> 0.
  \quad \blacksquare\ \text{[Q.E.D.]}
\end{aligned}
\end{equation*}

{\raggedright
This property is verified in the
\href{https://github.com/thiagomata/prime-numbers/blob/master/src/main/scala/v1/chapter4/cycle/integral/recursive/properties/CycleIntegralProperties.scala}{\texttt{CycleIntegralProperties::\allowbreak assertCycleIntegralPositive}}.
\par}

\subsection{Unit-Cycle Generation of Consecutive Integers}

The simplest non-empty cycle is the unit cycle $[1]$: every step adds exactly
one. Its cycle integral enumerates the consecutive integers starting after
the initial value --- the candidate stream that the prime sieve filters in
later articles. Because each step contributes exactly one, the integral has
an exact closed form at every position, and strict increase follows from it
as the unit-cycle instance of Subsection~4.4.

\begin{equation*}
init \geq 0 \implies \operatorname{CycleIntegral}([1], init)_i = init + i + 1
\end{equation*}

\textbf{Proof.} Induction on $i$.

\textbf{Base Case} ($i = 0$):

\begin{equation*}
\begin{aligned}
\operatorname{Cycle}([1])_0 &= 1 && \text{[Unit Cycle]} \\
CI_0 &= init + \operatorname{Cycle}([1])_0 && \text{[By Definition]} \\
&= init + 1 && \text{[Substitution]}
\end{aligned}
\end{equation*}

\textbf{Induction Step} ($i > 0$):

\begin{equation*}
\begin{aligned}
CI_{i-1} &= init + (i-1) + 1 && \text{[Induction Hypothesis]} \\
&= init + i && \text{[Simplification]} \\
\operatorname{Cycle}([1])_i &= 1 && \text{[Unit Cycle]} \\
CI_i &= CI_{i-1} + \operatorname{Cycle}([1])_i && \text{[Step Property, Subsection~3.1]} \\
&= init + i + 1 && \text{[Substitution]}
\end{aligned}
\end{equation*}

\begin{equation*}
\therefore\ \forall\, i \in \mathbb{N}_0:\ \operatorname{CycleIntegral}([1], init)_i = init + i + 1
\quad \blacksquare\ \text{[Q.E.D.]}
\end{equation*}

{\raggedright
This property is verified in the
\href{https://github.com/thiagomata/prime-numbers/blob/master/src/main/scala/v1/chapter4/cycle/integral/recursive/properties/CycleIntegralOnesProperties.scala}{\texttt{CycleIntegralOnesProperties::\allowbreak assert\allowbreak Cycle\allowbreak Integral\allowbreak Of\allowbreak Ones}}.
A key Scala verification excerpt is in Appendix~A.17; the complete proof is
linked in the source reference.
\par}

\textbf{Strict increase.} For positions $0 \leq a < b$, the closed form
gives

\begin{equation*}
\begin{aligned}
CI_b - CI_a &= (init + b + 1) - (init + a + 1) && \text{[Unit-Cycle Closed Form]} \\
&= b - a && \text{[Simplification]} \\
&> 0 && \text{[Since } b > a\text{]}
\end{aligned}
\end{equation*}

\begin{equation*}
\therefore\ 0 \leq a < b \implies \operatorname{CycleIntegral}([1], init)_b > \operatorname{CycleIntegral}([1], init)_a
\quad \blacksquare\ \text{[Q.E.D.]}
\end{equation*}

{\raggedright
This property is verified in the
\href{https://github.com/thiagomata/prime-numbers/blob/master/src/main/scala/v1/chapter4/cycle/integral/recursive/properties/CycleIntegralOnesProperties.scala}{\texttt{CycleIntegralOnesProperties::\allowbreak assert\allowbreak Cycle\allowbreak Integral\allowbreak Of\allowbreak Ones\allowbreak Strictly\allowbreak Increasing}}.
A key Scala verification excerpt is in Appendix~A.17.
\par}
